\section{Exact shuffle quotients and a two-product proof at weight five}
\label{sec:shuffle}

\subsection{The quotient being computed}

Let \(\mathcal H^1\) be the rational vector space on the empty word
and all binary words ending in \(1\), with shuffle multiplication.
It is bigraded by word length \(w\) and number of ones \(d\).
Let \(\mathcal J\) be its augmentation ideal, spanned by nonempty
words. The indecomposable quotient
\[
Q_{w,d}=(\mathcal J/\mathcal J^2)_{w,d}
\]
sets products of two positive-weight elements equal to zero.
This is a formal combinatorial construction. Shuffle equalities evaluate
to analytic identities at \(z\); the corresponding map on indecomposable
quotients requires also quotienting the evaluated algebra by products.
Evaluation can introduce further relations. In particular, the formal
dimension is not the dimension of a space of numerical periods.

\begin{theorem}[The exact formal shuffle count]\label{thm:witt}
For \(w\ge1\) and \(1\le d\le w\),
\[
\dim_{\Q}Q_{w,d}
=\ell_{w,d}
:=\frac1w\sum_{k\mid\gcd(w,d)}
\mu(k)\binom{w/k}{d/k}.
\]
A basis is represented by the binary Lyndon words of length \(w\)
with \(d\) ones. Every such word belongs to \(\mathcal H^1\).
Every other word has a finite, exact polynomial expression in the full
family of Lyndon generators, each monomial having the original total
bidegree. The number of independent product relations in
this component is
\[
\binom{w-1}{d-1}-\ell_{w,d}.
\]
\end{theorem}

\begin{proof}
Order the alphabet by \(0<1\) and words lexicographically.
A Lyndon word is strictly smaller than every nontrivial rotation.
Every word has a unique Chen--Fox--Lyndon factorization into a
nonincreasing concatenation of Lyndon words.
In characteristic zero, shuffle monomials in Lyndon words form a
polynomial basis \cite{Radford}. One can see the constructive point
directly: if the factorization is \(u_1^{m_1}\cdots u_r^{m_r}\) with
distinct \(u_1>\cdots>u_r\), then the shuffle product
\(u_1^{\shuffle m_1}\shuffle\cdots\shuffle u_r^{\shuffle m_r}\)
has this concatenation as its lexicographically largest word, with
coefficient \(m_1!\cdots m_r!\). All other words are smaller.
Induction therefore gives a triangular polynomial expression, and the
same triangularity proves linear independence of the monomials.

Every Lyndon word of length greater than one that contains both letters
starts with \(0\) and ends with \(1\). The single-letter words are
\(0\) and \(1\). The factorization of a word ending in \(1\) cannot
contain the factor \(0\), because that smallest factor would be last,
forcing the concatenation to end in \(0\).
Thus \(\mathcal H^1\) is exactly the polynomial algebra on all Lyndon
words other than \(0\). Passing to its augmentation quotient leaves
one class for each generator.

It remains to count. A primitive binary word of length \(w\) with
\(d\) ones has \(w\) distinct rotations and exactly one Lyndon rotation.
Every binary word is a power of a unique primitive word.
M\"obius inversion on this power decomposition gives
\(\sum_{k\mid\gcd(w,d)}\mu(k)\binom{w/k}{d/k}\)
primitive words, proving the formula. Finally, there are
\(\binom{w-1}{d-1}\) words ending in \(1\) in the component, which
gives the claimed product rank.
\end{proof}

\begin{table}[htbp]
\centering
\begin{tabular}{c|rrrrr}
\toprule
Weight \(w\)&\(d=1\)&\(d=2\)&\(d=3\)&\(d=4\)&\(d=5\)\\
\midrule
3&1&1&0&--&--\\
4&1&1&1&0&--\\
5&1&2&2&1&0\\
6&1&2&3&2&1\\
7&1&3&5&5&3\\
8&1&3&7&8&7\\
9&1&4&9&14&14\\
\bottomrule
\end{tabular}
\caption{The exact formal indecomposable counts \(\ell_{w,d}\).
They count Lyndon generators, not independent numerical constants.}
\label{tab:witt}
\end{table}

For \(w=6,d=3\), the ten words reduce modulo products to three Lyndon
words:
\[
000111,\qquad001011,\qquad001101,
\]
corresponding respectively to \(F_{4,1,1}\), \(F_{3,2,1}\),
and \(F_{3,1,2}\). This explains the manuscript's finite rank-seven
observation and gives its uniform extension.
The companion code records the full polynomial normal forms of all ten
triples. For a compact view of the quotient, put
\(X=[F_{4,1,1}],Y=[F_{3,2,1}],Z=[F_{3,1,2}]\), where brackets mean
classes modulo products. Exact triangular elimination gives
\[
\begin{array}{c|rrr@{\qquad}c|rrr}
 &X&Y&Z&&X&Y&Z\\ \hline
{}[F_{4,1,1}]&1&0&0 &[F_{1,4,1}]&6&3&1\\
{}[F_{3,2,1}]&0&1&0 &[F_{1,3,2}]&-12&-6&-2\\
{}[F_{3,1,2}]&0&0&1 &[F_{1,2,3}]&12&6&2\\
{}[F_{2,3,1}]&-9&-4&-1 &[F_{1,1,4}]&-6&-3&-1\\
{}[F_{2,2,2}]&12&4&0 &[F_{2,1,3}]&-3&-1&0
\end{array}
\]
For example, one full identity before taking the quotient is
\begin{equation}
F_{2,2,2}(z)=12F_{4,1,1}(z)+4F_{3,2,1}(z)
-2\Li_2(z)F_{3,1}(z)+\frac{\Li_2(z)^3}{6}.
\label{eq:full-222}
\end{equation}
Re-expanding its right side into words gives the single word \(010101\)
with coefficient one. This supplies a finite rational certificate,
including the product terms that a quotient table suppresses.

\subsection{An elementary all-weight double elimination}

At depth two the count simplifies to
\[
\ell_{w,2}=\left\lfloor\frac{w-1}{2}\right\rfloor.
\]
There is also a direct triangular proof that avoids general shuffle
algebra. For \(p+q=w\), shuffling the two single-polylogarithm words gives
\begin{align}
\Li_p(z)\Li_q(z)
={}&\sum_{j=1}^{p}\binom{w-j-1}{q-1}F_{w-j,j}(z)
+\sum_{j=1}^{q}\binom{w-j-1}{p-1}F_{w-j,j}(z).
\label{eq:general-double-shuffle}
\end{align}
The binomial coefficient counts the interleavings of the zero blocks
before the second nonzero letter.

\begin{proposition}[Independent double product rows]\label{prop:double-rank}
For fixed \(w\), the rows of
\eqref{eq:general-double-shuffle} with
\(1\le p\le\lfloor w/2\rfloor\) have rank
\(\lfloor w/2\rfloor\) on the \(w-1\) double coordinates.
They express the entire family through products of singles and the
\(\lfloor(w-1)/2\rfloor\) coordinates
\[
\{F_{w-j,j}:1\le j\le\lfloor(w-1)/2\rfloor\}.
\]
This set is empty when \(w=2\).
\end{proposition}

\begin{proof}
In the row indexed by \(p\), the largest inner index appearing is
\(q=w-p\). Its coefficient is \(1\) when \(p<q\), and \(2\)
when \(p=q\). No larger inner index occurs.
Order the rows with \(p\) decreasing and use the columns with
\(j=w-p\); the corresponding square block is triangular with nonzero
diagonal. All products of two depth-one words have already been included,
since swapping \(p,q\) repeats a row. This proves the rank and the
elimination statement.
\end{proof}

At \(z=i\), taking imaginary parts of these equations gives an
unconditional spanning bound for the \(g_{a,b}\), with explicit
single-value terms. It can be sharpened in the parity sector, as the
next section shows.

\subsection{The weight-five relation is two shuffles}

The words for \(\Li_1\Li_4\) and \(\Li_2\Li_3\) give
\begin{align}
\Li_1(z)\Li_4(z)
&=2F_{4,1}(z)+F_{3,2}(z)+F_{2,3}(z)+F_{1,4}(z),
\label{eq:shuffle14}\\
\Li_2(z)\Li_3(z)
&=6F_{4,1}(z)+3F_{3,2}(z)+F_{2,3}(z).
\label{eq:shuffle23}
\end{align}
The exact single values are
\[
\Li_1(i)=-\frac L2+\frac{\pi i}{4},\quad
\Li_2(i)=-\frac{\pi^2}{48}+iG,\quad
\Li_3(i)=-\frac{3\zeta(3)}{32}+\frac{\pi^3i}{32},
\]
\[
\Li_4(i)=-\frac{7\pi^4}{11520}+i\beta(4).
\]
Consequently
\begin{align}
2g_{41}+g_{32}+g_{23}+g_{14}
&=-\frac12\beta(4)L-\frac{7\pi^5}{46080},
\label{eq:wt5-first}\\
6g_{41}+3g_{32}+g_{23}
&=-\frac{\pi^5}{1536}-\frac{3G\zeta(3)}{32}.
\label{eq:wt5-second}
\end{align}
Here \(g_{41}\) means \(g_{4,1}\), and similarly for other subscripts.

\begin{corollary}[Proof and strengthening of the recorded relation]
\label{cor:wt5}
The manuscript's identity
\[
576g_{41}+288g_{32}+736g_{23}+960g_{14}
=21G\zeta(3)-480\beta(4)L
\]
holds. Moreover the four \(g\)-values are spanned by \(g_{41},g_{32}\)
and the single-value products \(\pi^5,G\zeta(3),\beta(4)L\).
Explicitly,
\begin{align}
g_{23}
&=-6g_{41}-3g_{32}-\frac{\pi^5}{1536}
  -\frac{3G\zeta(3)}{32},\label{eq:g23-elim}\\
g_{14}
&=4g_{41}+2g_{32}+\frac{23\pi^5}{46080}
  +\frac{3G\zeta(3)}{32}-\frac12\beta(4)L.
\label{eq:g14-elim}
\end{align}
\end{corollary}

\begin{proof}
Multiply \eqref{eq:wt5-first} by \(960\), multiply
\eqref{eq:wt5-second} by \(-224\), and add.
The \(\pi^5\) coefficients cancel exactly. Solving the two
independent equations instead gives \eqref{eq:g23-elim} and
\eqref{eq:g14-elim}.
\end{proof}

This proof explains why a search can return an apparently complicated
relation while missing the simpler structure: eliminating a chosen
single-value coordinate can enlarge coefficients.
The mathematical content is visible before the elimination.

\subsection{The surviving mixed relation}

For the harmonic sums used by the manuscript,
\[
S_p=\sum_{n\ge0}\frac{(-1)^nH_n}{(2n+1)^p},
\]
the exact bridge to colored doubles is
\begin{equation}
S_p=\Im\bigl(F_{p,1}(i)+\Li_{p,1}(i,-1)\bigr).
\label{eq:S-mixed}
\end{equation}
Indeed, for an odd outer index \(m=2n+1\), the inner coefficient
\(\sum_{k<m}(1+(-1)^k)/k\) is \(H_n\); even outer indices contribute
zero to the imaginary part. Absolute convergence justifies the argument
for \(p>1\), and Abel limits cover \(p=1\).

The recorded \(S_4\) candidate becomes, after the proved
\eqref{eq:g23-elim}, the more economical statement
\begin{conjecture}[A remaining mixed Gaussian reduction]\label{conj:S4}
\begin{equation}
S_4=\frac{58g_{41}+24g_{32}}7
       +\frac{19\pi^5}{3584}-2\beta(4)\log2.
\label{eq:S4-short}
\end{equation}
\end{conjecture}
This is equivalent to the existing candidate, not a newly discovered
independent relation. Independent integral evaluations support it,
but the present article does not supply its functional-equation proof.
In particular, it must not be silently added to the exact shuffle
certificate. It identifies a concrete additional colored relation to
seek. Neither its numerical support nor the dimensions of a selected
finite relation system establish arithmetic independence of any
surviving coordinate.
